\subsection{Decisiones del diseño según Jain}
\label{sec:jain_diseno}
El enfoque de \citet{jain1991} y su material docente \citep{jain2017} orientan la definición del sistema, la carga, las respuestas, los factores y los controles. El servicio del sistema es ejecutar un grafo y entregar sus resultados declarados. Se usan cinco patrones con archivos de tamaños e identidades conocidos. En distribución se distinguen particiones diferentes y consumidores de un mismo archivo. Las respuestas principales son los volúmenes en bytes, DME y los conteos. La duración del flujo y los tiempos posteriores se miden por separado.

Cada contraste tiene una línea base y varía el tamaño, la agrupación o la ubicación según la pregunta. Se conserva el resto de los controles pertinentes y se verifica su valor efectivo con el plan y el historial. Los niveles de las fichas explican qué cambia y por qué. Los contrastes permiten comprobar el cálculo por archivo y destino. Un diseño factorial permitiría examinar interacciones entre los factores con niveles comparables en todos los patrones. El requerido se deriva de la definición; las cuentas del observado se reconstruyeron retrospectivamente desde contratos preservados y no constituyen predicciones previas al experimento.

\subsection*{Cobertura de la selección}
La Tabla~\ref{tab:fuera_contrastes} identifica las nueve condiciones elegibles que quedaron fuera de las fichas. La selección retrospectiva permite comprobar las cuentas de cada contraste. Estimar la frecuencia de estos comportamientos en otros flujos requeriría otro diseño. Las otras condiciones documentan asignaciones alternativas, cambios en el número de tareas, un nivel adicional de agrupación o cambios de pool y ubicación. Sus 27 ejecuciones se conservan como archivo de campaña y no integran las tablas de DME de esta tesis. El cotejo detallado de 348 direcciones de este capítulo corresponde a las 45 ejecuciones seleccionadas.

\begin{table}[htbp]\centering\small
\caption{Condiciones elegibles fuera de las fichas de validación}
\label{tab:fuera_contrastes}
\begin{tabularx}{\linewidth}{@{}P{3.3cm}Y@{}}
\toprule Condiciones & Qué añaden al bloque A\\\midrule
L02, D01 & Asignaciones alternativas del pipeline y de distribución.\\
L05 & Pipeline más largo, de cinco etapas.\\
A02B, D02B, R02B & Grafos de agregación, distribución y redistribución con menos ramas.\\
GD2 & Nivel intermedio entre D02 y GD4 para agrupar particiones distintas.\\
W01, W04 & Cambios conjuntos de asignación y número de máquinas disponibles.\\
\bottomrule
\end{tabularx}
\end{table}
